\documentclass[10pt]{article}
\usepackage{../../framework/SI_Interros}
\usepackage{array}
\setlength{\parindent}{0pt}
\renewcommand{\arraystretch}{1.15}
\newcommand{\mobility}[6]{\begin{tabular}{c|c}R&T\\\hline #1&#2\\#3&#4\\#5&#6\end{tabular}}
\begin{document}
\InterroHeader{Modélisation des mécanismes}{20}
\question{Donner la relation qui relie le pas $p$ d'un dispositif vis--écrou, la position $z(t)$ et la position angulaire $\beta(t)$.}{1}
\begin{reponse}[22mm]\end{reponse}
\question{Dans chacun des cas suivants, compléter les informations manquantes : symbole cinématique normalisé, nom de la liaison et tableau de mobilités. Utiliser deux couleurs différentes pour les symboles cinématiques normalisés.}{3}
\begin{minipage}[t]{.78\linewidth}\centering
\begin{tabular}{|>{\raggedright\arraybackslash}p{.27\linewidth}|>{\centering\arraybackslash}p{.53\linewidth}|>{\centering\arraybackslash}p{.13\linewidth}|}\hline
\textbf{Désignation de la liaison} & \textbf{Symbole plan} & \textbf{Tableau de mobilités}\\\hline
Hélicoïdale d'axe $(C,\vec y)$ & \rule{0pt}{17mm} & \\\hline
\rule{0pt}{17mm} & $+\,H$ & \mobility{0}{1}{0}{0}{0}{0}\\\hline
Pivot d'axe $(B,\vec x)$ & \rule{0pt}{17mm} & \\\hline
\end{tabular}\end{minipage}\hfill
\begin{minipage}[t]{.18\linewidth}\centering\vspace{20mm}
\begin{tikzpicture}[>=Latex,scale=1.0]\draw[->,very thick](0,0)--(1.7,0)node[right]{$x$};\draw[->,very thick](0,0)--(0,1.7)node[above]{$y$};\draw[->,very thick](0,0)--(.35,-.35)node[below right]{$z$};\fill(0,0)circle(2pt);\end{tikzpicture}\end{minipage}
\question{Combien y a-t-il de fonctions AGIR dans l'attacheur de vignes ?}{1}\begin{reponse}[13mm]\end{reponse}
\question{Quel est le moteur utilisé sur la barrière Sympact ?}{1}\begin{reponse}[13mm]\end{reponse}
\question{Quel constituant du pilote hydraulique permet de convertir l'énergie mécanique de rotation en énergie hydraulique ?}{1}\begin{reponse}[13mm]\end{reponse}
\newpage
\InterroHeader{Modélisation des mécanismes}{20}
\par\medskip
Soit la modélisation d'une pompe à pistons axiaux :
\begin{center}
\begin{tikzpicture}[scale=.9,>=Latex]
\draw[thick] (-2.6,-1.2)--(-1.4,1.4);\foreach \y in {-1.0,-.7,...,1.2}{\draw (-2.65,\y)--(-2.35,\y+.2);} 
\draw[very thick,blue] (-1.7,.45)--(1.3,.45);\draw[blue,fill=white,very thick](-1.7,.45)circle(.17) node[above right,black]{$A$};
\draw[red,thick,fill=white](-.1,.2)rectangle(.6,.7) node[midway,black]{$B$};\draw[red,thick](.25,.2)--(.25,-.65)--(.05,-.65);
\draw[red,thick,fill=white](.05,-.85)rectangle(.9,-.45) node[above right,black]{$C$};\draw[blue,very thick](.6,.45)--(2.0,.45) node[right,black]{2};\draw[red,very thick](.9,-.65)--(2.0,-.65) node[right,black]{$x_0$};
\draw[->](-1.25,-.6)--(-.3,-.05)node[right]{$\vec x$};\draw[->](-1.25,-.6)--(-1.25,.9)node[above]{$\vec y_0$};\node at (.1,-.25){1};
\end{tikzpicture}
\end{center}
\question{Donner le graphe des liaisons correspondant.}{2}\begin{reponse}[67mm]\end{reponse}
\question{Donner les torseurs cinématiques $\{\mathcal V_{2/1}\}_B$, $\{\mathcal V_{1/0}\}_C$ et $\{\mathcal V_{2/0}\}_A$ associés aux trois liaisons du schéma cinématique.}{2}\begin{reponse}[105mm]\end{reponse}
\end{document}
